\chapter{What These Two Domains Together Establish}
\label{ch:two-domains}

\begin{plainlanguage}
\textbf{In plain terms:} Two unrelated case studies produced the same underlying pattern — a tempting coincidence this chapter deliberately refuses to over-interpret. Two hand-picked cases are not statistical proof of anything general; they're a lead worth following, stated honestly as that and nothing more.
\end{plainlanguage}

\section{The structure that recurred}

Two domains, chosen for having essentially nothing else in common --- sovereign flood-control infrastructure and a toddler's vocabulary --- produced the same qualitative pattern when the same naive ordering/threshold hypothesis was tested against each:
\[
\mathcal{S} = (\Delta t,\ N/O)
\]
--- a joint structure combining measurable interface lag between coupled subsystems and a separate nominal-versus-operational closure distinction. Both domains showed: (a) at least one pair of subsystems with near-zero lag, consistent with high coupling; (b) at least one subsystem whose triggering event or originating timescale was independent of the others, breaking the assumption of a shared reference point; and (c) a distinction between declared/apparent closure and actual underlying reorganization, with the gap between them sometimes actively maintained by a mechanism that benefits from the appearance of resolution (the levee suppressing the signal for restructuring; production-level overextension persisting despite intact comprehension underneath it).

\section{What this recurrence does and does not establish}

It is tempting, having found the same structure twice, to claim the structure is general --- a shared mechanism operating across institutional and cognitive-developmental systems alike. That claim should be resisted in exactly the form it is tempting to state it.

\begin{observation}[Structural Resemblance, not ``Cross-Domain Invariance'']
The same formal pattern $\mathcal{S}$ was observed in two domains chosen for reasons unrelated to testing this book's formalism --- one institutional, one developmental --- and was not the pattern either domain was originally selected to illustrate.
\end{observation}

\subsection*{Why similar structures do not imply shared causes}

It is worth stating explicitly, and not leaving to inference, exactly what kind of mistake the demotion above is designed to prevent --- because the temptation it guards against is easy to fall into without noticing. Finding the same pattern in two domains licenses, at most, the inference \emph{same pattern}; it does not license \emph{same mechanism}. These are different claims, and the gap between them is not a technicality. Two systems can exhibit identical formal structure for entirely unrelated underlying reasons --- a predator-prey population cycle and an electrical oscillator both solve the same class of differential equation, without either being a hidden cause of the other, or evidence that biology and electronics share a common substrate. Structural resemblance is consistent with shared mechanism, consistent with coincidence, and consistent with a deeper, currently unidentified common cause distinct from either domain studied --- and nothing in finding the resemblance twice distinguishes among these three possibilities. What would distinguish them is a mechanism-level account connecting the two domains directly, which this book has not attempted and does not claim to have found. The honest content of this chapter's finding is exactly as stated above: a recurring formal pattern, motivating further investigation, supplying useful shared vocabulary for later chapters --- and nothing about \emph{why} the pattern recurs.

A natural-sounding but unjustified next step would assert something like $P(\mathcal{S}\mid D_1, D_2) > P(\mathcal{S}\mid D_1)\cdot P(\mathcal{S}\mid D_2)$ --- that finding the pattern in both domains is itself evidence the domains share an underlying mechanism, by analogy to how correlated observations across independent samples support a common cause. This inference does not go through here. It presupposes a sampling process --- domains drawn at random from some larger population of domains, with $\mathcal{S}$ then checked for co-occurrence --- that does not exist in this book. Two domains were chosen, not sampled; and they were not selected entirely independently of each other, since the institutional case was investigated first and the developmental case was investigated partly because earlier work in this project had already developed vocabulary (subsystem indexing, interface lag) that made the developmental literature's production/comprehension dissociation easy to recognize once encountered. With $n=2$ non-randomly-selected, non-independently-selected cases, no probability statement of this form is estimable, and stating one would borrow a statistical authority this evidence has not earned --- precisely the error Chapter~\ref{ch:fidelity}'s machinery exists to catch, here applied to the book's own argument rather than to an external case.

What can honestly be said: the same distinction recurred in two domains selected for unrelated reasons, which motivates further investigation into whether it is a general feature of coupled repair systems, and supplies the book's working vocabulary (interface lag, coupling, nominal/operational closure) for the chapters that follow. It does not establish a shared mechanism. A reader entitled to remain unconvinced that anything general has been shown here is not making an error; they are correctly applying the standard this book has set for itself.

\section{What carries forward}

Three pieces of machinery, built and tested rather than merely asserted, carry into the rest of the book: the interface-lag concept and its open coupling conjecture; the nominal/operational closure distinction, which will recur explicitly in Part~VI's discussion of borrowed authority, where a declared elimination and an operationally real one turn out to be exactly the same distinction in a different domain; and the anti-vacuity principle governing when a global process may legitimately be split into subsystems --- a principle Part~IV will need immediately, the first time this book asks whether ``intelligence,'' like ``the child's category graph,'' might also turn out to be more than one thing wearing a single name.
